\scanpage{7}
qui est un torseur \inserted{à droite} \unsure{sous|sur} $\unsure{\mathfrak S|\mathscr S}_l$. Les op. $\sigma_0,\ldots,\sigma_l$ sur $R\times\operatorname{Rep}(I)$ s'obtiennent manifestement ainsi : $\sigma_0,\sigma_1,\ldots,\sigma_{l-1}$ sont les opérations $\operatorname{id}_R\times\rho_{\tau_{i+1}}$ ($0\leq i\leq l-2$) où $\tau_{i+1}\in\mathfrak S_l$ est la transposition de $i+1$ et $i+2$ ; \inserted{i.e.}
\[
\sigma_i(r,u:[1,l]\overset{\sim}{-}I)=(r,u\circ\tau_{\unsure{i|i'}})\qquad 0\leq i\leq l-1
\]
quant à $\sigma_{l-1}$, on obtient
\[
\sigma_{l-1}(r,u:[1,l]\xrightarrow{\sim}I)=\{\varphi(r,u(l))\mathbin{.}r,u\}
\]
On trouve bien que \inserted{$\sigma_i^2=1$ ($0\leq i\leq l-1$) et que} $\sigma_i$ $\sigma_j$ commutent si $j\geq i+2$ ($0\leq i,j\leq l-1$) -- le seul cas non \unsure{tt à ft|tout à ft} trivial étant $j=l-1$, $i\leq l-3$,
\[
\begin{aligned}
\sigma_i\sigma_{l-1}(r,u)&=\sigma_i(\varphi(r,u(l)).r,u)=(\varphi(r,u(l)).r,u\circ\tau_{i+1})\\
\sigma_{l-1}\sigma_i(r,u)&=\sigma_{l-1}(r,u\circ\tau_{i+1})
=\varphi(r,\underbrace{(u\circ\tau_{i+1})(l)}_{\substack{\Vert\\u(l)\ \text{si}\\i\leq l-3\ \text{i.e.}\\i+1\leq l-2}}).r,u\circ\tau_{i+1})
\end{aligned}
\]
Ici $(\sigma_i\sigma_{i+1})^3=\operatorname{id}$ -- distinguer d'abord le cas $0\leq i\leq l-2$ -- où celà résulte de $(\tau_{i+1}\tau_{i+2})^3=1$, puis $i=l-1$, où $\sigma_{l-1}\sigma_{l-2}$ est donné par
\[
\begin{aligned}
\sigma_{l-1}\sigma_{l-2}(r,u)&=\sigma_{l-1}(r,u\circ\tau_{l-1})=\varphi(r,(u\circ\tau_{l-1})(l)).r,\unsure{u\circ\tau_{l-1}}\\
\pi(r,u)&=(\underbrace{\overbrace{\varphi(r,u(l-1))}^{w}.r}_{r'},\underbrace{u\circ\tau_{l-1}}_{u'})\\
\pi^2(r,u)&=\inserted{\pi(r',u')=}\deleted{\unsure{\varphi(\varphi(r,u(l-1)).r,\ ,u|\varphi(\varphi(r,u(l-1))r,u)}}\\
&(\underbrace{\varphi(r',\underbrace{u'(l-1)}_{u(l)}).r'}_{\varphi(wr,u(l))\,w\,r,\,u},\underbrace{u'\circ\tau_{l-1}}_{u})
\end{aligned}
\]
% archive sheet number: 6

